%%%%%%%%%%%%%%%%%%%%%%%%%%%%%%%%%%%%%%%%%%%%%%%%%%%%%%%%%%%%%%%%%%%%%%%%%%%%%%%%
% almamater Beamer style / Template, one build per preset
% Andrea Omicini
% Alma Mater Studiorum - Università di Bologna
% mailto:andrea.omicini@unibo.it
%
% This template is released into the public domain under CC0 1.0 Universal:
% use it, adapt it and build on it freely, with no obligation to credit.
% See the LICENSE file, or https://creativecommons.org/publicdomain/zero/1.0/
%
% The almamater style it uses is covered by its own licence (LPPL 1.3c):
% https://github.com/andreaomicini/beamer-almamater
%%%%%%%%%%%%%%%%%%%%%%%%%%%%%%%%%%%%%%%%%%%%%%%%%%%%%%%%%%%%%%%%%%%%%%%%%%%%%%%%
% the preset: AMSUniBo unless the build names another (the CI builds all six, -usepretex)
\providecommand{\almamaterpreset}{AMSUniBo}
\input{presets/\almamaterpreset/options}
\documentclass[presentation,\almamaterseal,apice]{beamer}\mode<presentation>{\usetheme{\almamaterpreset}}
%%%%%%%%%%%%%%%%%%%%%%%%%%%%%%%%%%%%%%%%%%%%%%%%%%%%%%%%%%%%%%%%%%%%%%%%%%%%%%%%
\def\BibTeX{{\rm B\kern-.05em{\sc i\kern-.025em b}\kern-.08em
    T\kern-.1667em\lower.7ex\hbox{E}\kern-.125emX}}
\renewcommand{\thefootnote}{\fnsymbol{footnote}}
\input{presets/\almamaterpreset/preamble}
% version -- Major.Minor[.Patch]; the release time-stamp is appended by the CI
\newcommand{\templatemajor}{2}
\newcommand{\templateminor}{0}
% \templatepatch is optional: comment the next line out to release as Major.Minor
\newcommand{\templatepatch}{0}
\ifdefined\templatepatch{}
  \newcommand{\templateversion}{\templatemajor.\templateminor.\templatepatch}
\else
  \newcommand{\templateversion}{\templatemajor.\templateminor}
\fi
%%%%%%%%%%%%%%%%%%%%%%%%%%%%%%%%%%%%%%%%%%%%%%%%%%%%%%%%%%%%%%%%%%%%%%%%%%%%%%%%
%
\subtitle[AMS Series Templates]
{AMS Series Templates}
%
\author[\sspeaker{Omicini} \and Else]
{\speaker{Andrea Omicini} \and Someone Else}
%
\institute[DISI, Univ. Bologna]
{\url{mailto:andrea.omicini@unibo.it} \qquad \url{mailto:someone.else@unibo.it}\\
Dipartimento di Informatica \textendash{} Scienza e Ingegneria (DISI)\\
\textsc{Alma Mater Studiorum} \textendash{} Università di Bologna}
%
\date[v.\ \templateversion]{version \templateversion}
%
%%%%%%%%%%%%%%%%%%%%%%%%%%%%%%%%%%%%%%%%%%%%%%%%%%%%%%%%%%%%%%%%%%%%%%%%%%%%%%%%
\begin{document}
%%%%%%%%%%%%%%%%%%%%%%%%%%%%%%%%%%%%%%%%%%%%%%%%%%%%%%%%%%%%%%%%%%%%%%%%%%%%%%%%

%/////////
\frame{\titlepage}
%/////////

%%===============================================================================
%\section*{Outline}
%%===============================================================================
%
%%/////////
%\frame[c]{\tableofcontents[hideallsubsections]}
%%/////////

%===============================================================================
\section{Premises}
%===============================================================================

%/////////
\begin{frame}[c]{Academic Presentations}
%
\begin{itemize}
	\item technical presentations are common in the academic environment
	\item \alert{\LaTeX{}}\footnote{\uurl{https://www.latex-project.org}} is typically used for scientific literature
	\begin{itemize}
		\item for better formulas, bibliography, typography, \ldots
	\end{itemize}
	\item so, \LaTeX{} is the most typical \emph{toolbox} for scientific presentations
	\begin{itemize}
		\item which easily translates into teaching presentations, too
	\end{itemize}
\end{itemize}
%
\end{frame}
%/////////

%/////////
\begin{frame}[c]{Beamer}
%
\begin{itemize}
	\item \alert{Beamer} is a \LaTeX{} document class for creating presentation \emph{slides}, with a wide range of \emph{templates} and a set of features for making \emph{slideshow effects}
	\begin{description}
		\item[repo] \uurl{https://github.com/josephwright/beamer}
		\item[CTAN] \uurl{https://ctan.org/pkg/beamer}
	\end{description}
	\item even though other styles for \LaTeX{} presentations exist, Beamer is by far the best way today
\end{itemize}
%
\end{frame}
%/////////

%/////////
\input{presets/\almamaterpreset/ams}
%%
%/////////

%/////////
\input{presets/\almamaterpreset/aim}
%%
%/////////

%/////////
\begin{frame}[c]{Who}
%
\begin{itemize}
	\item mostly, myself and my collaborators
	\item everybody else, be my guest
\end{itemize}
%
\end{frame}
%/////////

%/////////
\begin{frame}[c]{Distribution}
%
\begin{block}{Public on GitHub}
	\begin{itemize}
		\item the style is distributed on \uurl{https://github.com/andreaomicini/beamer-almamater/}
		\item the releases are available at \uurl{https://github.com/andreaomicini/beamer-almamater/releases}
		\item this template is distributed on \uurl{https://github.com/andreaomicini/beamer-almamater-template/}
	\end{itemize}
\end{block}
%
\end{frame}
%/////////

%===============================================================================
\section{Use}
%===============================================================================

%/////////
\begin{frame}[c,fragile]{Declaration}
%
\begin{block}{\texttt{\textbackslash{}documentclass}}
Your main Beamer \texttt{.tex} file should open with the declaration
%
% set in type rather than verbatim, so as to name this build's preset
\par\smallskip
{\ttfamily\hspace*{2em}\textbackslash documentclass[presentation]\{beamer\}\par
\hspace*{4em}\textbackslash mode<presentation>\{\textbackslash usetheme\{\almamaterpreset\}\}\par}
\smallskip
%
so as to use the \almamaterpreset{} preset of the almamater Beamer style
\end{block}
%
\end{frame}
%/////////

%/////////
\begin{frame}[c,fragile]{Bibliography Style}
%
\begin{block}{\texttt{apalike-AMS}}
Your main Beamer \texttt{.tex} file should include the declaration
\begin{verbatim}
    \bibliographystyle{apalike-AMS}
\end{verbatim}
so as to use the almamater \BibTeX{} style
\end{block}
%
\end{frame}
%/////////

%/////////
\begin{frame}[c,fragile]{Template}
%
\begin{block}{\texttt{almamater-template.tex}}
This template's sources can be used as a simple example of how to use this Beamer style
\end{block}
%
\end{frame}
%/////////


%===============================================================================
\section{Style}
%===============================================================================

%/////////
\input{presets/\almamaterpreset/colours}
%%
%/////////

%/////////
\input{presets/\almamaterpreset/palette}
%%
%/////////

%/////////
\begin{frame}[c,fragile]{Blocks}
%
\begin{block}{This is a \texttt{block} environment}
\begin{verbatim}
\begin{block}{...}
...
\end{block}
\end{verbatim}
\end{block}
%
\begin{exampleblock}{This is an \texttt{exampleblock} environment}
\begin{verbatim}
\begin{exampleblock}{...}
...
\end{exampleblock}
\end{verbatim}
\end{exampleblock}
%
\begin{alertblock}{This is an \texttt{alertblock} environment}
\begin{verbatim}
\begin{alertblock}{...}
...
\end{alertblock}
\end{verbatim}
\end{alertblock}
%
\end{frame}
%/////////

%/////////
\input{presets/\almamaterpreset/surfaces}
%%
%/////////

%/////////
\begin{frame}[c,fragile]{One Command, Every Surface}
%
The same source line, here and in each header and body below:\\
\emph{emph}, \alert{alert}, \aalert{aalert}\ccite{bibtex-patashnik88}, \uurl{apice.unibo.it}, \uuurl{apice.unibo.it}
%
\begin{block}{\emph{emph}, \alert{alert}, \aalert{aalert}\ccite{bibtex-patashnik88}, \uurl{apice.unibo.it}, \uuurl{apice.unibo.it}}
\emph{emph}, \alert{alert}, \aalert{aalert}\ccite{bibtex-patashnik88}, \uurl{apice.unibo.it}, \uuurl{apice.unibo.it}
\end{block}
%
\begin{alertblock}{\emph{emph}, \alert{alert}, \aalert{aalert}\ccite{bibtex-patashnik88}, \uurl{apice.unibo.it}, \uuurl{apice.unibo.it}}
\emph{emph}, \alert{alert}, \aalert{aalert}\ccite{bibtex-patashnik88}, \uurl{apice.unibo.it}, \uuurl{apice.unibo.it}
\end{alertblock}
%
\begin{exampleblock}{\emph{emph}, \alert{alert}, \aalert{aalert}\ccite{bibtex-patashnik88}, \uurl{apice.unibo.it}, \uuurl{apice.unibo.it}}
\emph{emph}, \alert{alert}, \aalert{aalert}\ccite{bibtex-patashnik88}, \uurl{apice.unibo.it}, \uuurl{apice.unibo.it}
\end{exampleblock}
%
{\tiny\begin{verbatim}
\emph{emph}, \alert{alert}, \aalert{aalert}\ccite{bibtex-patashnik88},
\uurl{apice.unibo.it}, \uuurl{apice.unibo.it}
\end{verbatim}} %chktex 31
%
\end{frame}
%/////////

%===============================================================================
\section{New Commands}
%===============================================================================

%/////////
\input{presets/\almamaterpreset/citations}
%%
%/////////

%/////////
\input{presets/\almamaterpreset/urls}
%%
%/////////

%/////////
\input{presets/\almamaterpreset/doi}
%%
%/////////

%/////////
\input{presets/\almamaterpreset/apice-addresses}
%%
%/////////

%/////////
\input{presets/\almamaterpreset/alert-emphasis}
%%
%/////////

%/////////
\input{presets/\almamaterpreset/speakers}
%%
%/////////

%===============================================================================
\section{New Options}
%===============================================================================

%/////////
\begin{frame}[c,fragile,allowframebreaks]{\texttt{apice} option}
%
\begin{alertblock}{\texttt{apice} Beamer option}
\begin{itemize}
	\item the option \texttt{apice} can be used in the presentation preamble
	\begin{itemize}
		\item as in
		\begin{verbatim}
\documentclass[presentation,apice]{beamer}...
		\end{verbatim}
	\end{itemize}
	\item enables the display of the \BibTeX{} field \texttt{apice}
	\begin{itemize}
		\item containing the relative URL of the publication on the \alert{\href{https://apice.unibo.it/}{\textsf{APICe}}} XWiki
	\end{itemize}
	\item properly linked to the corresponding \href{https://apice.unibo.it/}{\textsf{APICe}} page
\end{itemize}
\end{alertblock}
%
\begin{exampleblock}{Example}
%
The following \BibTeX{} entry in the \texttt{.bib} file with the \texttt{apice} option enabled
%
{\tiny\begin{verbatim}
@manual{bibtex-patashnik88,
    apice = {BibtexPatashnik88},
    author = {Patashnik, Oren},
    booktitle = {\BibTeXing},
    organization = {CTAN, The Comprehensive TeX Archive Network},
    title = {\BibTeX{}ing},
    url = {https://mirrors.ctan.org/biblio/bibtex/base/btxdoc.pdf},
    year = 1988}
\end{verbatim}} %chktex 31
%
would translate into the following entry in the presentation bibliography
%
\begin{thebibliography}{}
\bibitem[Patashnik, 1988]{bibtex-patashnik1988}
Patashnik, O. (1988)
\newblock {\bfseries {\em \BibTeX{}ing}}
\newblock CTAN, The Comprehensive TeX Archive Network
\newblock\apicepubpar{BibtexPatashnik88} \uuurl{https://mirrors.ctan.org/biblio/bibtex/base/btxdoc.pdf}
\end{thebibliography}
%
\end{exampleblock}
%
\end{frame}
%/////////

%/////////
\input{presets/\almamaterpreset/seal-options}
%%
%/////////

%===============================================================================
\section{Bibliography}
%===============================================================================

%/////////
\begin{frame}[c]{Entry Types}
%
\begin{block}{What \texttt{apalike-AMS} adds to \texttt{apalike}}
\begin{itemize}
	\item every classical type\textemdash{}\texttt{@article}, \texttt{@book}, \texttt{@inproceedings}, \ldots{}\textemdash{}as in \texttt{apalike}, titles in bold
	\item \texttt{@talk}\textemdash{}a talk, as APICe's Talks write it\ccite{impossibleconsensus-slides2016}
	\item \texttt{@issue}\textemdash{}a whole journal issue, by its editors\ccite{agentlink3-jaamas17}
	\item theses\textemdash{}\texttt{@phdthesis}, \texttt{@mastersthesis} with APICe's fields: supervisors, cycle, AMS e-print\ccite{sita-thesis2017}
	\item \texttt{subtitle}\textemdash{}after the title, past a colon\ccite{agentsandmasdevelopment-book2026}
	\item \texttt{articleno}, \texttt{numpages}\textemdash{}an article number in place of pages\ccite{lptech4mas-jaamas35}
	\item \texttt{@webpage}\textemdash{}a web page, every field but a link optional\ccite{beamer-ctan,transformer-googleblog2017}
	\item last line of every entry\textemdash{}the APICe marker, then the DOI, else the URL
\end{itemize}
\end{block}
%
\end{frame}
%/////////

%/////////
\begin{frame}[c,fragile]{\texttt{@talk}}
%
\begin{block}{A talk, as APICe's Talks write it\ccite{impossibleconsensus-slides2016}}
{\scriptsize\begin{verbatim}
@talk{impossibleconsensus-slides2016,
    address = {London, UK},
    apice = {ImpossibleconsensusQcon2016},
    author = {Howard, Heidi},
    date = {2016-03-29},
    howpublished = {QCon London 2016},
    speaker = {Howard, Heidi},
    title = {Distributed Consensus: Making Impossible Possible},
    type = {Talk},
    year = 2016}
\end{verbatim}} %chktex 31
\begin{itemize}
	\item the venue line\textemdash{}\texttt{type}, \texttt{howpublished}, \texttt{address}, \texttt{date} written out
	\item \texttt{speaker} only when it differs from \texttt{author}
\end{itemize}
\end{block}
%
\end{frame}
%/////////

%/////////
\begin{frame}[c,fragile]{\texttt{@issue}}
%
\begin{block}{A whole journal issue, by its editors\ccite{agentlink3-jaamas17}}
{\scriptsize\begin{verbatim}
@issue{agentlink3-jaamas17,
    apice = {Agentlink3Jaamas17},
    editor = {Omicini, Andrea and McBurney, Peter},
    journal = {Autonomous Agents and Multi-Agent Systems},
    number = 3,
    title = {{S}pecial Issue on Foundations, Advanced Topics
             and Industrial Perspectives of Multi-agent Systems},
    volume = 17,
    year = 2008}
\end{verbatim}} %chktex 31
\begin{itemize}
	\item editors as in \texttt{@proceedings}, journal and volume as in \texttt{@article}
\end{itemize}
\end{block}
%
\end{frame}
%/////////

%/////////
\begin{frame}[c,fragile]{Theses}
%
\begin{block}{A thesis, with APICe's fields\ccite{sita-thesis2017}}
{\scriptsize\begin{verbatim}
@mastersthesis{sita-thesis2017,
    amseprint = {13847},
    apice = {SitaLogicmiddleware2017},
    author = {Sita, Alberto},
    cosupervisor = {Denti, Enrico and Calegari, Roberta},
    cycle = {LM},
    supervisor = {Omicini, Andrea},
    title = {Agenti, programmazione logica e sistemi distribuiti:
             esperimenti in JADE e tuProlog},
    year = 2017}
\end{verbatim}} %chktex 31
\begin{itemize}
	\item supervisors printed; with neither DOI nor URL, the AMS e-print links the thesis
\end{itemize}
\end{block}
%
\end{frame}
%/////////

%/////////
\begin{frame}[c,fragile]{Subtitles, Article Numbers}
%
\begin{block}{\texttt{subtitle}\ccite{agentsandmasdevelopment-book2026}}
{\scriptsize\begin{verbatim}
    subtitle = {Platforms, Toolkits, Technologies},
    title = {Agents and Multi-Agent Systems Development},
\end{verbatim}} %chktex 31
\begin{itemize}
	\item the subtitle follows the title, past a colon
\end{itemize}
\end{block}
%
\begin{block}{\texttt{articleno}, \texttt{numpages}\ccite{lptech4mas-jaamas35}}
{\scriptsize\begin{verbatim}
    articleno = 1,
    numpages = 67,
    pages = {1--67},
\end{verbatim}} %chktex 31
\begin{itemize}
	\item \texttt{articleno} replaces \texttt{pages}, and \texttt{numpages} follows it\textemdash{}\emph{article 1, 67 pages}
\end{itemize}
\end{block}
%
\end{frame}
%/////////

%/////////
\begin{frame}[c,fragile]{\texttt{@webpage}}
%
\begin{block}{A web page\textemdash{}also \texttt{@online}, \texttt{@www}, \texttt{@electronic}}
{\tiny\begin{verbatim}
@webpage{beamer-ctan,
    title = {beamer},
    url = {https://ctan.org/pkg/beamer}}

@webpage{transformer-googleblog2017,
    apice = {TransformerGoogleblog2017},
    author = {Uszkoreit, Jakob},
    organization = {Google Research},
    title = {Transformer: A Novel Neural Network Architecture
             for Language Understanding},
    url = {https://ai.googleblog.com/2017/08/transformer-novel-neural-network.html},
    year = 2017}
\end{verbatim}} %chktex 31
\end{block}
%
\end{frame}
%/////////

%/////////
\begin{frame}[c]{\texttt{@webpage}: Who, Title}
%
\begin{block}{What comes first}
\begin{itemize}
	\item every field but a link optional\textemdash{}nothing to invent for a bare page\ccite{beamer-ctan}
	\item first the author, else \texttt{organization}, with the year\ccite{transformer-googleblog2017}
	\item with neither, the title opens the entry
	\item the title in bold, \alert{as written}\textemdash{}a web page's title is usually a name
\end{itemize}
\end{block}
%
\end{frame}
%/////////

%/////////
\begin{frame}[c]{\texttt{@webpage}: Links, Labels, Warnings}
%
\begin{block}{How a web page prints}
\begin{itemize}
	\item the URL on its own line, in the colour of a journal's name
	\item APICe and DOI last, as for every type\textemdash{}a DOI alone is a link enough
	\item the label is the author's, else \texttt{key}, else the organisation, else the title\textemdash{}\texttt{[beamer]}
	\begin{itemize}
		\item \texttt{, <year>} only when there is a year; two equal labels with none read \texttt{[Kubernetes, a]}, \texttt{[Kubernetes, b]}
	\end{itemize}
	\item \texttt{urldate} is not printed, as in every type
	\item a warning only with neither \texttt{url} nor \texttt{doi}
\end{itemize}
\end{block}
%
\begin{alertblock}{No \texttt{@url}}
\texttt{url} is a field, so \BibTeX{} cannot have a type of that name: an \texttt{@url} entry prints as \texttt{@misc}\textemdash{}write \texttt{@webpage}
\end{alertblock}
%
\end{frame}
%/////////

%===============================================================================
\section*{} 
%===============================================================================

%/////////
\frame{\titlepage}
%/////////

%===============================================================================
\section*{\refname}
%===============================================================================

%%%%
\setbeamertemplate{page number in head/foot}{}
%/////////
%\begin{frame}[c,noframenumbering]{\refname}
\begin{frame}[t,allowframebreaks,noframenumbering]{\refname}
	\tiny
%	\scriptsize
%	\footnotesize
	\bibliographystyle{apalike-AMS}
	\bibliography{almamater-template}
\end{frame}
%/////////

%%%%%%%%%%%%%%%%%%%%%%%%%%%%%%%%%%%%%%%%%%%%%%%%%%%%%%%%%%%%%%%%%%%%%%%%%%%%%%%%
\end{document}
%%%%%%%%%%%%%%%%%%%%%%%%%%%%%%%%%%%%%%%%%%%%%%%%%%%%%%%%%%%%%%%%%%%%%%%%%%%%%%%%
